% !TEX TS-program = pdflatexmk
\documentclass[12pt]{amsart}

%\usepackage[parfill]{parskip}    % Activate to begin paragraphs with an empty line rather than an indent

\usepackage[margin=1in]{geometry}

\usepackage{amsmath,amssymb,amsthm,latexsym,graphicx}
\usepackage[normalem]{ulem}
\usepackage{setspace} %used for doublespacing, etc.
\usepackage{hyperref}
\usepackage[shortlabels]{enumitem}
\usepackage{cancel}
\usepackage[dvipsnames,usenames]{color}
\usepackage[all]{xy}
\DeclareGraphicsRule{.tif}{png}{.png}{`convert #1 `dirname #1`/`basename #1 .tif`.png}

%Some useful environments.
\newtheorem{theorem}{Theorem}
\newtheorem{corollary}[theorem]{Corollary}
\newtheorem{conjecture}[theorem]{Conjecture}
\newtheorem{lemma}[theorem]{Lemma}
\newtheorem{proposition}[theorem]{Proposition}
\newtheorem{definition}[theorem]{Definition}
\newtheorem{exercise}[theorem]{Exercise}
\newtheorem{axiom}{Axiom}
\theoremstyle{remark}
\newtheorem{remark}{Remark}

%Some useful shortcuts for our favorite fields
\newcommand{\RR}{{\mathbb R}}
\newcommand{\NN}{{\mathbb N}}
\newcommand{\ZZ}{{\mathbb Z}}
\newcommand{\QQ}{{\mathbb Q}}
\newcommand{\CC}{{\mathbb C}}
\newcommand{\FF}{{\mathbb F}}


%Some useful shortcuts for formatting lists
\newcommand{\bc}{\begin{center}}
\newcommand{\ec}{\end{center}}
\newcommand{\be}{\begin{enumerate}}
\newcommand{\ee}{\end{enumerate}}
\newcommand{\bi}{\begin{itemize}}
\newcommand{\ei}{\end{itemize}}

%Some useful shortcuts for formatting mathematical symbols
\newcommand{\ol}[1]{\overline{#1}}
\newcommand{\oimp}[1]{\overset{#1}{\Longleftrightarrow}} %labeled iff symbol
\newcommand{\bv}[1]{\ensuremath{ \vec{\mathbf{#1}}} } %makes a vector.
\newcommand{\mc}[1]{\ensuremath{\mathcal{#1}}} %put something in caligraphic font
\newcommand{\mpg}[1]{\marginpar{ #1}} %to put comments in margins
\newcommand{\bsl}[1]{\texttt{\symbol{92}{\em #1}}} %for backslashes.
\newcommand{\normale}{\trianglelefteq}
\newcommand{\normal}{\triangleleft}

%Commenting tools
\usepackage{soul}
\definecolor{highlight}{rgb}{1,0.6,0.6}
\sethlcolor{highlight}
\newcommand{\hlm}[1]{\colorbox{highlight}{$\displaystyle #1$}}
\newtheoremstyle{mycomment}{\topsep}{-0in}{\small \itshape \sffamily}{}{\small \itshape\sffamily}{:}{.5em}{}
\theoremstyle{mycomment}
\newtheorem*{acomment}{\color{BrickRed}{Comment}}
\newcommand{\com}[1]{{\color{OliveGreen}\begin{acomment}{#1} %#2 \color{black} 
\end{acomment}\noindent}}
%\newcommand{\com}[1]{{\color{BrickRed}{\\ Comment:}\color{OliveGreen}{#1} \\}}
\newcommand{\red}[1]{{\color{BrickRed} #1}}
\def\midd{\,\middle\vert\, }
\newcommand{\blue}[1]{{\color{MidnightBlue}#1}}
\newcommand{\green}[1]{{\color{OliveGreen}#1}}
\newcommand{\mwrong}[2]{\red{\cancel{#1}}\green{#2}}
\newcommand{\wrong}[2]{\red{\sout{#1}}\green{#2}}
\definecolor{OliveGreen}{rgb}{.3,.5,.2}
\definecolor{MidnightBlue}{rgb}{.3,.4,.6}

\title{Assignment \#12 \& \#13 \today}
\DeclareMathOperator{\lcm}{lcm}

\author{Coordinator: Stefano Fochesatto}
\begin{document}
\maketitle
%\setstretch{2.5} %Use for 2.5 spacing
\doublespacing %Use for double spacing


% 10.1.11, 
% 10.2.4, 
% 10.2.5, 
% 11.1.1
% 11.1.2,

\section*{Section 10.1}


\begin{exercise}[\textbf{10.1.11}] Let $M$ be the abelian group $\ZZ/24\ZZ \times \ZZ/15\ZZ \times \ZZ/50\ZZ$.
  \begin{enumerate}
    \item[(a)] Find the annihilator of $M$ in $\ZZ$ (i.e, a generator for this principal ideal).
    \begin{proof} Consider the ideal generated by the $l = \lcm(24, 15, 50)$. We will proceed by showing that $(l) = ann_{\ZZ}(M)$.
      Suppose $a \in (l)$ so $a = lj$ some $j$ multiple of $l$ in $\ZZ$
      and let $m \in M$ such that $m = (a, b, c)$ where $a \in \ZZ/24\ZZ$, $b \in \ZZ/15\ZZ$ and $c \in \ZZ/50\ZZ$. 
      Then note that since $l = \lcm(24, 15, 50)$ we get $am = (lja, ljb, ljc) = (0, 0, 0)$. So $a \in ann_{\ZZ}(M)$ and $(l)\subseteq ann_{\ZZ}(M)$. 
      
      Let $j \in ann_{\ZZ}(M)$. Then by definition for all $m = (a, b, c) \in M$ we get that $jm = (ja, jb, jc) = 0$. So it must follow that
      that $j(1, 1, 1) = (j, j, j) = (0, 0, 0)$ so it must be the case that $j \in (\lcm(24, 15, 50)) = (l)$. Therefore
      $ann_{\ZZ}(M) \subseteq (l)$, thus finally $(l) = ann_{\ZZ}(M)$
    \end{proof}
\vspace{.15in}

    \item[(b)] Let $I = 2\ZZ$. Describe the annihilator of $I$ in $M$ as a direct product of cyclic groups. 
   \begin{proof}

    \end{proof}
\vspace{.15in}
  \end{enumerate}

\end{exercise}
\vspace{.5in}



\section*{Section 10.2}



\begin{exercise}[\textbf{10.2.4}] Let $A$ be any $\ZZ$-module, let $a$ be any element of $A$ and let $n$ be a positive integer. 
  Prove that the map $\varphi_a: \ZZ/n\ZZ \to A$ given by $\varphi_a(\overline{k}) = ka$ is a well defined $\ZZ$-module homomorphism if and only if $na = 0$
  Prove that $Hom_\ZZ(\ZZ/n\ZZ, A)\cong A_n$, where $A_n = \{ a \in A |na = 0\}$.    
    \begin{proof} Let $A$ be any $\ZZ$-module, let $a$ be any element of $A$ and let $n$ be a positive integer. Suppose the map  $\varphi_a: \ZZ/n\ZZ \to A$ given by $\varphi_a(\overline{k}) = ka$ is a well defined $\ZZ$-module homomorphism. Note that $na = \varphi_a(\overline{n}) = \varphi_a(\overline{0}) = 0a = a$. Now suppose $na = 0$ and let $i \cong k \mod n$, so by definition we know that there exists some 
      $j \in \NN$ such that $i = nj + k$. Note that $\varphi_a(\overline{i}) = ia = (nj + k)a = nja + ka = j(na) + ka = ka = \varphi_a(\overline{k})$. Therefore $\varphi_a$ is well-defined.





      Consider the mapping $\pi: A_n \to Hom_\ZZ(\ZZ/n\ZZ, A)$ defined by $\pi(a) = \varphi_a$. Suppose $\varphi_i, \varphi_j \in Hom_\ZZ(\ZZ/n\ZZ, A)$ such that $\varphi_i = \varphi_j$, then it follows that 
      $\varphi_i(\overline{1}) = 1i = i$ and $\varphi_j(\overline{1}) = 1j = j$ so $i = j$. Therefore this map is one to one. Let $\varphi_n \in Hom_\ZZ(\ZZ/n\ZZ, A)$ note that since $A$ is a $\ZZ$-module 
    


    
  \end{proof}

\end{exercise}
\vspace{.5in}




\begin{exercise}[\textbf{10.2.5}] Exhibit all $\ZZ$-module homomorphisms from $\ZZ/30\ZZ$ to $\ZZ/21\ZZ$.
  \begin{proof}
    
  \end{proof}

\end{exercise}
\vspace{.5in}



\section*{Section 11.1}

\begin{exercise}[\textbf{11.1.1}] Let $V = R^n$ and let $(a_1, a_2, \dots, a_n)$ be a fixed vector in $V$. Prove that the collection of elements
  $(z_1, x_2, \dots, x_n)$ of $V$ with $a_1x_1 + a_2x_2 + a_3x_3 + \dots + a_nx_n = 0$ is a subspace of $V$. Determine the dimension of this subspace 
  and find a basis.  
  \begin{proof} 
    
  \end{proof}
\end{exercise}
\vspace{.5in}



\begin{exercise}[\textbf{11.1.2}] Let $V$ be the collection of polynomials with coefficients in $\QQ$ in the variable $x$ of degree
  at most 5. Prove that $V$ is a vector space over $\QQ$ of dimension 6, with $1,x,x^2, \dots, x^5$ as a basis. Prove that $1, 1 + x, 1 + x + x^2, \dots, 1 + x + x^2 + x^3 + x^4 + x^5$ is 
  also a basis for $V$.
  \begin{proof}
    
  \end{proof}
\end{exercise}
\vspace{.5in}







\end{document}